\begin{itemize}
\item \textbf{Ionizační energie:} energie nutná k odtržení nejslaběji vázaného elektronu, $\mathrm A\to\mathrm{A^+}+e^-$. První ionizační energie se vztahuje k neutrálnímu atomu, druhá k odtržení dalšího elektronu z kationtu.
\item \textbf{Elektronová afinita:} energie uvolněná při přijetí volného elektronu neutrálním atomem, $\mathrm A+e^-\to\mathrm{A^-}$.
\item \textbf{Elektronegativita:} schopnost atomu přitahovat elektrony v chemické vazbě.
\end{itemize}
Ionizační energie a afinita se udávají v eV na částici nebo kJ/mol; elektronegativita je bezrozměrná.
